% !TeX program = pdflatex
\documentclass[aspectratio=169,xcolor={dvipsnames,table}]{beamer}

\input{../../_en_preambulo_beamer}
% Comandos y macros estadísticas compartidas para las presentaciones Beamer en inglés (Metropolis)

% Conjuntos numéricos básicos
\newcommand{\R}{\mathbb{R}}
\newcommand{\N}{\mathbb{N}}
\newcommand{\Q}{\mathbb{Q}}
\newcommand{\Z}{\mathbb{Z}}
\newcommand{\set}[1]{\left\{ #1 \right\}}
\newcommand{\abs}[1]{\left|#1\right|}
\newcommand{\norm}[1]{\left\| #1 \right\|}

% Comandos estadísticos y probabilísticos del curso
\newcommand{\Var}{\operatorname{Var}}
\newcommand{\cov}{\operatorname{Cov}}
\newcommand{\corr}{\rho}
\newcommand{\comb}[2]{\begin{pmatrix} #1 \\ #2 \end{pmatrix}}
\newcommand{\s}{\sigma}
\newcommand{\particion}{\mathfrak{P}}
\newcommand{\card}[1]{\#\left( #1 \right)}
\newcommand{\seq}[3]{\set{#1}_{#2}^{#3}}
\newcommand{\E}{\mathbb{E}}
\newcommand{\Prob}{\mathbb{P}}

% Entornos pedagógicos y cajas en inglés académico natural (soportando tanto nombres en inglés como en español)
\newenvironment{discussion}{%
  \begin{alertblock}{Discussion Question}%
}{%
  \end{alertblock}%
}
\newenvironment{discusion}{%
  \begin{alertblock}{Discussion Question}%
}{%
  \end{alertblock}%
}

\newenvironment{reflection}{%
  \begin{alertblock}{Classroom Reflection}%
}{%
  \end{alertblock}%
}
\newenvironment{reflexion}{%
  \begin{alertblock}{Classroom Reflection}%
}{%
  \end{alertblock}%
}

\newenvironment{paradox}{%
  \begin{alertblock}{Exploratory Paradox}%
}{%
  \end{alertblock}%
}
\newenvironment{paradoja}{%
  \begin{alertblock}{Exploratory Paradox}%
}{%
  \end{alertblock}%
}

\newenvironment{motivation}{%
  \begin{exampleblock}{Motivation: Data Science in Practice}%
}{%
  \end{exampleblock}%
}
\newenvironment{motivacion}{%
  \begin{exampleblock}{Motivation: Data Science in Practice}%
}{%
  \end{exampleblock}%
}

\newenvironment{quickchallenge}{%
  \begin{block}{Quick Team Challenge}%
}{%
  \end{block}%
}
\newenvironment{retoclase}{%
  \begin{block}{Quick Team Challenge}%
}{%
  \end{block}%
}


\title{Expectation and Variance}
\subtitle{Section 03.03 --- Moments of Discrete Variables}
\author[J. Castillo Colmenares]{Juliho Castillo Colmenares}
\institute[Tec de Monterrey]{Tecnológico de Monterrey}
\date{\vspace{-1.2cm}}

\begin{document}

% SLIDE 01: Cover Page
\begin{frame}[plain]
	\titlepage
\end{frame}

% SLIDE 02: Roadmap
\begin{frame}{Roadmap --- Chapter 03: Discrete Random Variables}
	\begin{block}{Curricular Structure of Chapter 03}
		\small
		\begin{enumerate}
			\itemsep=0.04cm
			\item \textcolor{gray}{Section 03.01: Discrete Probability Mass Functions (PMF and Support)}
			\item \textcolor{gray}{Section 03.02: Cumulative Distribution Functions (Discrete CDF)}
			\item \textbf{\textcolor{TecRojo}{Section 03.03: Mathematical Expectation, Variance, and Moments}} $\leftarrow$ \emph{Current focus}
			\item \textcolor{gray}{Section 03.04: Bernoulli and Binomial Distributions}
			\item \textcolor{gray}{Section 03.05: Geometric and Negative Binomial Distributions}
			\item \textcolor{gray}{Section 03.06: Hypergeometric Distribution}
		\end{enumerate}
	\end{block}
	\vspace{-0.15cm}
	\begin{alertblock}{Curricular Objective of the Section}
		\footnotesize
		Master the algebraic foundations and operational properties of expected value and variance, applying the LOTUS Theorem and the probabilistic center of gravity theorem to decision-making and engineering problems.
	\end{alertblock}
\end{frame}

% SLIDE 03: Motivation and Intuition
\begin{frame}{Motivation --- From Distributions to Summary Numerical Summaries}
	\begin{columns}[T]
		\begin{column}{0.48\textwidth}
			\begin{block}{Complexity of Full PMFs}
				\small
				\begin{itemize}
					\itemsep=0.08cm
					\item A complete PMF $f_X(x)$ or CDF $F_X(x)$ contains the **full stochastic characterization** of a random variable $X$.
					\item In complex engineering systems or large-scale data pipelines, comparing entire distributions is algebraically intractable.
					\item We require **synthetic scalar invariants** that summarize central tendency and dispersion.
				\end{itemize}
			\end{block}
		\end{column}
		\begin{column}{0.48\textwidth}
			\begin{alertblock}{Center of Gravity and Inertia}
				\small
				\begin{itemize}
					\itemsep=0.08cm
					\item **Expectation ($\mu$):** The physical center of mass of the probability distribution, weighting every value by its likelihood.
					\item **Variance ($\sigma^2$):** The moment of inertia around $\mu$, quantifying risk, volatility, and quadratic spread.
				\end{itemize}
			\end{alertblock}
		\end{column}
	\end{columns}
\end{frame}

% SLIDE 04: Formal Definition of Expectation
\begin{frame}{Formal Definition --- Mathematical Expectation}
	\begin{block}{Analytical Definition of Population Mean $\mu = \mathbb{E}[X]$}
		\small
		Let $X$ be a discrete random variable with countable support $\mathcal{S}_X = \set{x_1, x_2, \dots}$ and probability mass function $f_X(x)$. The **mathematical expectation** or **expected value** of $X$ is defined as:
		\begin{align*}
			\mathbb{E}[X] = \sum_{x \in \mathcal{S}_X} x \cdot f_X(x),
		\end{align*}
		provided that the series converges absolutely ($\sum |x| f_X(x) < \infty$).
	\end{block}
	\vspace{-0.1cm}
	\begin{columns}[T]
		\begin{column}{0.48\textwidth}
			\begin{exampleblock}{Frequentist Interpretation}
				\footnotesize
				By the Law of Large Numbers, as independent trials are repeated $n \to \infty$ times, the sample arithmetic mean $\bar{X}_n$ converges almost surely to $\mathbb{E}[X]$.
			\end{exampleblock}
		\end{column}
		\begin{column}{0.48\textwidth}
			\begin{alertblock}{Conceptual Warning}
				\footnotesize
				$\mathbb{E}[X]$ does not necessarily belong to support $\mathcal{S}_X$. For instance, a fair six-sided die yields $\mathbb{E}[X] = 3.5 \notin \mathcal{S}_X$.
			\end{alertblock}
		\end{column}
	\end{columns}
\end{frame}

% SLIDE 05: LOTUS Theorem
\begin{frame}{LOTUS Theorem --- Law of the Unconscious Statistician}
	\begin{block}{Law of the Unconscious Statistician (LOTUS)}
		\small
		Let $X$ be a discrete random variable with PMF $f_X(x)$ and let $g: \mathbb{R} \to \mathbb{R}$ be any real-valued function. The expected value of the transformation $Y = g(X)$ is directly given by:
		\begin{align*}
			\mathbb{E}[g(X)] = \sum_{x \in \mathcal{S}_X} g(x) \cdot f_X(x).
		\end{align*}
	\end{block}
	\vspace{-0.15cm}
	\begin{columns}[T]
		\begin{column}{0.48\textwidth}
			\begin{exampleblock}{Absolute Operational Advantage}
				\footnotesize
				Computes $\mathbb{E}[g(X)]$ **directly from the original support $\mathcal{S}_X$ and PMF $f_X(x)$**, without needing to derive the induced PMF $f_Y(y)$ first.
			\end{exampleblock}
		\end{column}
		\begin{column}{0.48\textwidth}
			\begin{alertblock}{Algorithmic Efficiency}
				\footnotesize
				Bypasses the combinatorial overhead of grouping preimages $\set{x : g(x) = y}$ whenever $g(X)$ is non-injective or many-to-one.
			\end{alertblock}
		\end{column}
	\end{columns}
\end{frame}

% SLIDE 06: Properties of the Expectation Operator
\begin{frame}{Properties of the Expectation Operator}
	\begin{block}{Universal Linearity and Monotonicity}
		\small
		The expectation operator $\mathbb{E}[\cdot]$ is a **linear, additive, and homogeneous functional** over the space of random variables with finite mean:
		\begin{enumerate}
			\itemsep=0.06cm
			\item \textbf{Constant Invariance:} If $c \in \mathbb{R}$ is constant, then $\mathbb{E}[c] = c$.
			\item \textbf{Affine Linearity:} $\mathbb{E}[aX + b] = a\mathbb{E}[X] + b$ for all constants $a, b \in \mathbb{R}$.
			\item \textbf{Universal Additivity:} For any two variables $X$ and $Y$ (dependent or independent):
			\begin{align*}
				\mathbb{E}[aX + bY] = a\mathbb{E}[X] + b\mathbb{E}[Y].
			\end{align*}
			\item \textbf{Monotonicity:} If $X \ge Y$ almost surely, then $\mathbb{E}[X] \ge \mathbb{E}[Y]$.
		\end{enumerate}
	\end{block}
\end{frame}

% SLIDE 07: Operational Variance and König-Huygens
\begin{frame}{Operational Variance and König-Huygens Formula}
	\begin{columns}[T]
		\begin{column}{0.48\textwidth}
			\begin{block}{Core Definition of Variance}
				\small
				The variance of $X$, denoted $\text{Var}(X)$ or $\sigma_X^2$, is the second central moment around the mean:
				\begin{align*}
					\text{Var}(X) = \mathbb{E}\left[ (X - \mu)^2 \right].
				\end{align*}
				Measures average quadratic spread around $\mu$.
			\end{block}
		\end{column}
		\begin{column}{0.48\textwidth}
			\begin{alertblock}{Standard Deviation ($\sigma_X$)}
				\small
				The positive square root of variance:
				\begin{align*}
					\sigma_X = \sqrt{\text{Var}(X)}.
				\end{align*}
				Shares exact physical units with the random variable $X$.
			\end{alertblock}
		\end{column}
	\end{columns}
	\vspace{-0.1cm}
	\begin{exampleblock}{König-Huygens Identity (Computational Shortcut)}
		\small
		Expanding $(X-\mu)^2$ and applying linearity yields the fundamental operational formula:
		\begin{align*}
			\text{Var}(X) = \mathbb{E}[X^2] - (\mathbb{E}[X])^2.
		\end{align*}
	\end{exampleblock}
\end{frame}

% SLIDE 08: Raw and Central Moments
\begin{frame}{Taxonomy of Raw and Central Moments}
	\begin{block}{Hierarchy of Statistical Moments of Order $k \in \mathbb{Z}^+$}
		\scriptsize
		\centering
		\begin{tabular}{p{2.8cm}|p{2.8cm}|p{6.4cm}}
			\toprule
			\textbf{Moment Type} & \textbf{Formal Operator} & \textbf{Stochastic \& Engineering Meaning} \\
			\midrule
			Raw Moment ($k$) & $\alpha_k = \mathbb{E}[X^k]$ & Moments relative to origin ($x=0$). \\
			\midrule
			Central Moment ($k$) & $\mu_k = \mathbb{E}[(X - \mu)^k]$ & Moments centered at population mean. \\
			\midrule
			Skewness ($\gamma_1$) & $\gamma_1 = \mu_3 / \sigma^3$ & Tail asymmetry (left vs. right tails). \\
			\midrule
			Kurtosis ($\gamma_2$) & $\gamma_2 = (\mu_4 / \sigma^4) - 3$ & Tail heaviness and modal sharpness. \\
			\bottomrule
		\end{tabular}
	\end{block}
	\vspace{-0.15cm}
	\begin{alertblock}{Basic Properties of Central Moments}
		\footnotesize
		By definition, the first central moment is identically zero: $\mu_1 = \mathbb{E}[X - \mu] = 0$. The second central moment corresponds precisely to the variance: $\mu_2 = \text{Var}(X)$.
	\end{alertblock}
\end{frame}

% SLIDE 09: Standardization and Z-Scores
\begin{frame}{Standardization of Variables ($Z$-Scores)}
	\begin{block}{Affine Standardization Transformation}
		\small
		Given a discrete variable $X$ with mean $\mu_X$ and standard deviation $\sigma_X > 0$, the standardized random variable ($Z$-score) is defined as:
		\begin{align*}
			Z = \frac{X - \mu_X}{\sigma_X}.
		\end{align*}
	\end{block}
	\vspace{-0.15cm}
	\begin{columns}[T]
		\begin{column}{0.48\textwidth}
			\begin{exampleblock}{Universal Moment Invariance}
				\footnotesize
				By linearity of expectation and affine variance transformation rules:
				\begin{align*}
					\mathbb{E}[Z] = 0 \quad \text{and} \quad \text{Var}(Z) = 1.
				\end{align*}
			\end{exampleblock}
		\end{column}
		\begin{column}{0.48\textwidth}
			\begin{alertblock}{Practical Significance}
				\footnotesize
				Enables dimensionless comparisons across heterogeneous physical units, evaluating how many standard deviations an observation deviates from the mean.
			\end{alertblock}
		\end{column}
	\end{columns}
\end{frame}

% SLIDE 10: Center of Gravity Theorem
\begin{frame}{Probabilistic Center of Gravity Theorem}
	\begin{block}{Mean Squared Error Minimization}
		\small
		Let $X$ be a random variable with finite variance and let $a \in \mathbb{R}$ be an arbitrary scalar. The mean squared error centered at $a$ decomposes exactly as:
		\begin{align*}
			M(a) = \mathbb{E}\left[ (X - a)^2 \right] = \text{Var}(X) + (\mu - a)^2.
		\end{align*}
	\end{block}
	\vspace{-0.15cm}
	\begin{columns}[T]
		\begin{column}{0.48\textwidth}
			\begin{exampleblock}{Unique Minimum Solution}
				\footnotesize
				Since $(\mu - a)^2 \ge 0$, the absolute minimum is achieved if and only if $a = \mu$:
				\begin{align*}
					\min_{a \in \mathbb{R}} \mathbb{E}\left[ (X - a)^2 \right] = \text{Var}(X).
				\end{align*}
			\end{exampleblock}
		\end{column}
		\begin{column}{0.48\textwidth}
			\begin{alertblock}{Foundation of Least Squares}
				\footnotesize
				Establishes the population mean $\mu$ as the optimal point predictor under quadratic loss, where variance represents irreducible residual error.
			\end{alertblock}
		\end{column}
	\end{columns}
\end{frame}

% SLIDE 11A: Classroom Exercise --- Tier 1: Fundamental (Statement)
\begin{frame}{Classroom Exercise --- Tier 1: Fundamental (1/2)}
	\begin{block}{Problem 3.3.1 --- Error Transactions during Peak Hour in a Bank Database}
		\footnotesize
		Let $X$ be a discrete random variable modeling the number of error transactions during peak hour at a bank server, with support $\mathcal{S}_X = \{0, 1, 2, 3, 4\}$ and explicit mass function:
		\begin{table}[h]
			\centering
			\begin{tabular}{c|ccccc}
				\toprule
				$x$ & $0$ & $1$ & $2$ & $3$ & $4$ \\
				\midrule
				$f_X(x)$ & $0.35$ & $0.30$ & $0.20$ & $0.10$ & $0.05$ \\
				\bottomrule
			\end{tabular}
		\end{table}
		\begin{enumerate}\itemsep=0.1cm
			\item Formally compute the mathematical expectation $\mu = \mathbb{E}[X]$ using the fundamental analytical definition.
			\item Compute the second raw moment $\mathbb{E}[X^2]$ and determine the operational variance $\text{Var}(X)$ via the König-Huygens identity.
		\end{enumerate}
	\end{block}
\end{frame}

% SLIDE 11B: Classroom Exercise --- Tier 1: Fundamental (Solution)
\begin{frame}{Classroom Exercise --- Tier 1: Fundamental (2/2)}
	\fontsize{6.8pt}{8.2pt}\selectfont
	\begin{block}{1. Mathematical Expectation ($\mu = \mathbb{E}[X]$)}
		\fontsize{6.8pt}{8.2pt}\selectfont
		We weight each value of the variable by its population probability:
		$$\mu = \sum_{x=0}^4 x \cdot f_X(x) = 0(0.35) + 1(0.30) + 2(0.20) + 3(0.10) + 4(0.05)$$
		$$\mu = 0 + 0.30 + 0.40 + 0.30 + 0.20 = \mathbf{1.20\text{ transactions}}$$
	\end{block}

	\vspace{-0.15cm}
	\pause
	\begin{alertblock}{2. Second Raw Moment ($\mathbb{E}[X^2]$) and Operational Variance ($\text{Var}(X)$)}
		\fontsize{6.8pt}{8.2pt}\selectfont
		We compute the weighted average of the squares of the support:
		$$\mathbb{E}[X^2] = 0^2(0.35) + 1^2(0.30) + 2^2(0.20) + 3^2(0.10) + 4^2(0.05) = \mathbf{2.80}$$
		We apply the König-Huygens algebraic shortcut:
		$$\text{Var}(X) = \mathbb{E}[X^2] - \mu^2 = 2.80 - (1.20)^2 = 2.80 - 1.44 = \mathbf{1.36\text{ transactions}^2}$$
	\end{alertblock}
\end{frame}

% SLIDE 12A: Classroom Exercise --- Tier 2: Operational (Statement)
\begin{frame}{Classroom Exercise --- Tier 2: Operational (1/2)}
	\begin{block}{Problem 3.3.4 --- Mining Megaproject Bid and Risk-Adjusted Return}
		\footnotesize
		An engineering corporation evaluates a megaproject whose net profitability in millions of dollars ($R$) is conditioned by global scenarios:
		\begin{table}[h]
			\centering
			\begin{tabular}{l|ccc}
				\toprule
				\textbf{Economic Scenario} & \textbf{Probability ($p_k$)} & \textbf{Net Profitability ($r_k$)} \\
				\midrule
				Accelerated expansion & $0.25$ & $+120\text{ M USD}$ \\
				Moderate stability & $0.55$ & $+40\text{ M USD}$ \\
				Severe recession & $0.20$ & $-50\text{ M USD}$ \\
				\bottomrule
			\end{tabular}
		\end{table}
		\begin{enumerate}\itemsep=0.1cm
			\item Determine the expected net present value $\mathbb{E}[R]$ and the standard deviation of the return $\sigma_R$.
			\item If the board requires a financial safety metric $U = \mathbb{E}[R] - 1.25\sigma_R \ge 0$, should the project be accepted or rejected?
		\end{enumerate}
	\end{block}
\end{frame}

% SLIDE 12B: Classroom Exercise --- Tier 2: Operational (Solution)
\begin{frame}{Classroom Exercise --- Tier 2: Operational (2/2)}
	\fontsize{6.8pt}{8.2pt}\selectfont
	\begin{block}{1. Expected Value $\mathbb{E}[R]$ and Standard Deviation $\sigma_R$}
		\fontsize{6.8pt}{8.2pt}\selectfont
		$$\mathbb{E}[R] = 120(0.25) + 40(0.55) + (-50)(0.20) = 30.0 + 22.0 - 10.0 = \mathbf{+42.0\text{ M USD}}$$
		To obtain the risk $\sigma_R$, we first compute the second raw moment:
		$$\mathbb{E}[R^2] = (120)^2(0.25) + (40)^2(0.55) + (-50)^2(0.20) = 3600 + 880 + 500 = 4980$$
		$$\text{Var}(R) = \mathbb{E}[R^2] - (\mathbb{E}[R])^2 = 4980 - (42.0)^2 = 3216 \implies \sigma_R = \sqrt{3216} \approx \mathbf{56.71\text{ M USD}}$$
	\end{block}

	\vspace{-0.15cm}
	\pause
	\begin{alertblock}{2. Corporate Decision and Safety Threshold Evaluation $U$}
		\fontsize{6.8pt}{8.2pt}\selectfont
		We evaluate the risk-adjusted return indicator:
		$$U = \mathbb{E}[R] - 1.25\sigma_R = 42.0 - 1.25(56.71) = 42.0 - 70.89 = \mathbf{-28.89\text{ M USD}}$$
		Since the result is negative ($U < 0$), capital volatility exceeds the expected gain. The board \textbf{must formally reject} the bid due to excessive financial risk.
	\end{alertblock}
\end{frame}

% SLIDE 13A: Classroom Exercise --- Tier 3: Analytical (Statement)
\begin{frame}{Classroom Exercise --- Tier 3: Analytical (1/2)}
	\begin{block}{Problem 3.3.6 --- $Z$-Score Standardization and Invariant Moments}
		\footnotesize
		In an international civil engineering evaluation with score $X$, mean $\mu_X$, and standard deviation $\sigma_X > 0$, the standardized random variable ($Z$-score) is defined as:
		$$Z = \dfrac{X - \mu_X}{\sigma_X}.$$
		Formally prove, with full algebraic detail, that for any discrete random variable $X$, the standardized variable $Z$ rigorously satisfies the two universal identities:
		\begin{enumerate}\itemsep=0.15cm
			\item Zero mathematical expectation: $\mathbb{E}[Z] = 0$.
			\item Unit operational variance: $\text{Var}(Z) = 1$.
		\end{enumerate}
	\end{block}
\end{frame}

% SLIDE 13B: Classroom Exercise --- Tier 3: Analytical (Solution)
\begin{frame}{Classroom Exercise --- Tier 3: Analytical (2/2)}
	\fontsize{7.5pt}{9pt}\selectfont
	We can rewrite the transformation as a linear function $Z = cX + d$, where $c = \frac{1}{\sigma_X}$ and $d = -\frac{\mu_X}{\sigma_X}$. \pause

	\vspace{-0.05cm}
	\begin{block}{1. Proof of Zero Expectation ($\mathbb{E}[Z] = 0$)}
		\fontsize{7.5pt}{9pt}\selectfont
		Applying universal linearity ($\mathbb{E}[cX + d] = c\mathbb{E}[X] + d$) and recalling that $\mathbb{E}[X] = \mu_X$:
		$$\mathbb{E}[Z] = \mathbb{E}\left[ \dfrac{1}{\sigma_X} X - \dfrac{\mu_X}{\sigma_X} \right] = \dfrac{1}{\sigma_X}\mathbb{E}[X] - \dfrac{\mu_X}{\sigma_X} = \dfrac{\mu_X}{\sigma_X} - \dfrac{\mu_X}{\sigma_X} = \mathbf{0}. \quad \checkmark$$
	\end{block}

	\vspace{-0.1cm}
	\pause
	\begin{alertblock}{2. Proof of Unit Variance ($\text{Var}(Z) = 1$)}
		\fontsize{7.5pt}{9pt}\selectfont
		Applying the quadratic homogeneity property under scaling ($\text{Var}(cX + d) = c^2 \text{Var}(X)$) and recalling that the population variance of $X$ is $\text{Var}(X) = \sigma_X^2$:
		$$\text{Var}(Z) = \text{Var}\left( \dfrac{1}{\sigma_X} X - \dfrac{\mu_X}{\sigma_X} \right) = \left( \dfrac{1}{\sigma_X} \right)^2 \text{Var}(X) = \dfrac{\text{Var}(X)}{\sigma_X^2} = \dfrac{\sigma_X^2}{\sigma_X^2} = \mathbf{1}. \quad \blacksquare$$
	\end{alertblock}
\end{frame}

% SLIDE 14A: Classroom Exercise --- Tier 4: Challenging (Statement)
\begin{frame}{Classroom Exercise --- Tier 4: Challenging (1/2)}
	\begin{block}{Problem 3.3.9 --- St. Petersburg Paradox and Truncated Regularization}
		\footnotesize
		In the St. Petersburg Paradox, a fair coin is tossed repeatedly until the first heads appears on the $k$-th toss ($k \ge 1$), paying the player a reward of $X = 2^k$ coins.
		\begin{enumerate}\itemsep=0.12cm
			\item Prove that the theoretical expected value of the reward diverges to infinity ($\mathbb{E}[X] = \infty$).
			\item If, by capital regulation, the casino imposes a maximum bankroll reserve of $N$ tosses (if heads does not appear within the first $N$ attempts, the game ends paying $M = 2^N$), obtain a closed-form expression for the truncated expected value $\mathbb{E}[X_N]$ and evaluate it at $N = 25$ tosses.
		\end{enumerate}
	\end{block}
\end{frame}

% SLIDE 14B: Classroom Exercise --- Tier 4: Challenging (Solution)
\begin{frame}{Classroom Exercise --- Tier 4: Challenging (2/2)}
	\fontsize{7.5pt}{9pt}\selectfont
	\begin{block}{1. Theoretical Divergence of the Expectation $\mathbb{E}[X]$}
		\fontsize{7.5pt}{9pt}\selectfont
		The probability of obtaining the first heads on the $k$-th toss is $P(X=2^k) = \left(\frac{1}{2}\right)^k$:
		$$\mathbb{E}[X] = \sum_{k=1}^\infty 2^k P(X = 2^k) = \sum_{k=1}^\infty 2^k \left(\dfrac{1}{2}\right)^k = \sum_{k=1}^\infty 1 = 1 + 1 + 1 + \dots = \mathbf{\infty}$$
	\end{block}

	\vspace{-0.2cm}
	\pause
	\begin{alertblock}{2. Truncated Regularization via Bankroll Limit ($N=25$)}
		\fontsize{7.5pt}{9pt}\selectfont
		The truncated variable $X_N$ takes values $2^k$ for $k \in \{1, \dots, N\}$ with probability $(1/2)^k$, and the bounded value $2^N$ with residual probability $P(K > N) = (1/2)^N$:
		$$\mathbb{E}[X_N] = \sum_{k=1}^N 2^k \left(\dfrac{1}{2}\right)^k + 2^N \left(\dfrac{1}{2}\right)^N = \sum_{k=1}^N 1 + 1 = N + 1$$
		For a casino with a limit of $N = 25$ tosses ($\approx 33.5$ million coins in reserve), the fair entry price is extremely finite and modest:
		$$\mathbb{E}[X_{25}] = 25 + 1 = \mathbf{26\text{ exact coins}}.$$
	\end{alertblock}
\end{frame}

% SLIDE 18: Python Lab --- Raw and Central Moments
\begin{frame}[fragile]{Python Lab --- Moment Computation (`numpy`)}
	\begin{block}{Single Source of Truth Implementation (\texttt{03.03\_expectation\_and\_variance.py})}
		\lstinputlisting[language=Python, firstline=12, lastline=27, basicstyle=\fontsize{6.3pt}{7.5pt}\ttfamily]{../../code/03_variables_aleatorias_discretas/03.03_expectation_and_variance.py}
	\end{block}
\end{frame}

% SLIDE 19: Python Lab --- LOTUS and Z-Score Verification
\begin{frame}[fragile]{Python Lab --- LOTUS and $Z$-Score Verification}
	\begin{block}{Algebraic and Computational Invariance in `numpy`}
		\lstinputlisting[language=Python, firstline=32, lastline=49, basicstyle=\fontsize{6.3pt}{7.5pt}\ttfamily]{../../code/03_variables_aleatorias_discretas/03.03_expectation_and_variance.py}
	\end{block}
\end{frame}

% SLIDE 20: Conclusion & Monte Carlo Simulation of LLN
\begin{frame}[fragile]{Conclusion --- Monte Carlo Law of Large Numbers Simulation}
	\begin{columns}[T]
		\begin{column}{0.55\textwidth}
			\begin{block}{Sample Path Simulation ($N=100,000$)}
				\lstinputlisting[language=Python, firstline=54, lastline=66, basicstyle=\fontsize{5.6pt}{6.7pt}\ttfamily]{../../code/03_variables_aleatorias_discretas/03.03_expectation_and_variance.py}
			\end{block}
		\end{column}
		\begin{column}{0.42\textwidth}
			\begin{alertblock}{Key Conclusions}
				\scriptsize
				\begin{itemize}
					\itemsep=0.06cm
					\item $\bar{X}_n \to \mu = 1.20$ with tiny $0.0044$ error across $10^5$ iterations.
					\item $S_n^2 \to \sigma^2 = 1.36$ confirming ergodic empirical stability.
					\item Standardized $Z$-score variance is strictly unit ($1.000000$).
				\end{itemize}
			\end{alertblock}
		\end{column}
	\end{columns}
\end{frame}

\end{document}
